% GENERATED FILE. Edit canonical lessons, then run npm run generate:books.
% canonical-source-hash: fnv1a64:1a4098c6b96e0a38
% canonical-lessons: ZH-C393-zhufu, ZH-C393-qianwan, ZH-W393-wan-myriad, ZH-C393-wanneng, ZH-C393-suishi, ZH-C393-shunli

\chapter{Be Sure To}
\label{ch:zh-be-sure-to}
\hlchaptermodality{393}

\begin{chapteropening}
\textbf{By the end of this chapter:} \emph{I can talk about telling someone to remember, being sure to, what money can't do, any time and things going smoothly.}

Complete the last lesson of chapter 393: I can talk about telling someone to remember, being sure to, what money can't do, any time and things going smoothly.
\end{chapteropening}

\section[zhǔfu]{zhǔfu — to exhort; to tell (someone) to be sure to do something}

\label{lesson:ZH-C393-zhufu}



\begin{quote}
Before the new one: say the Mandarin for to change the name (of), then the Mandarin for to change one's tune.
\end{quote}

\subsection*{You'll want to know: zhǔfu}
\emph{zhǔfu} — \textquotedblleft{}to exhort; to tell (someone) to be sure to do something\textquotedblright{}.

\begin{quote}
\emph{māma zhǔfu wǒ wǎnshang zǎo huí jiā} — Mum told me to be sure to come home early tonight.
\end{quote}

Its characters come later; for now it is a word for the ear.

\begin{cousinweb}
\emph{zhǔfu} is built from \emph{zhǔ} 'to entrust' and \emph{fu}. The first is inherited from Middle Chinese \emph{tsyowk} (Guangyun, in Baxter's transcription), glossed 'to entrust to'; the dictionary has no entry for the second, so its history is not settled.
\end{cousinweb}

\begin{grammarlens}[title={What you've built}]
One more word for saying what would happen.
\end{grammarlens}

\subsection*{Your turn}
\begin{itemize}
\raggedright
  \item \emph{Say it:} \emph{zhǔfu}
  \item \emph{Say it:} \emph{zhǔfu}, once more
  \item \emph{Say it:} \emph{māma zhǔfu wǒ wǎnshang zǎo huí jiā} — Mum told me to be sure to come home early tonight.
  \item \emph{Recall:} say the Mandarin for to change the name (of), then the Mandarin for to change one's tune, then say \emph{zhǔfu} again
\end{itemize}

\begin{tcolorbox}[breakable,colback=teal!4,colframe=teal!35!black,title={Before you move on}]
What does zhǔfu mean? (\textquotedblleft{}To exhort; to tell (someone) to be sure to do something\textquotedblright{}.) Say it once more.
\end{tcolorbox}
\section[\texorpdfstring{\zh{千万} (qiānwàn)}{千万 (qiānwàn)}]{\zh{千万} (qiānwàn) — be sure to (urging); whatever you do}

\label{lesson:ZH-C393-qianwan}



\begin{quote}
Before the new one: say the Mandarin for to change one's tune, then the Mandarin for to exhort (someone).
\end{quote}

\subsection*{You'll want to know: \zh{千万}}
\textbf{\zh{千万}} — \emph{qiānwàn} — \textquotedblleft{}be sure to (urging); whatever you do\textquotedblright{}.

\begin{quote}
\textbf{\zh{你一个人在外边千万要保重}} \emph{nǐ yí ge rén zài wàibian qiānwàn yào bǎozhòng} — Out there on your own, be sure to look after yourself.
\end{quote}

The new character, in pieces: \zh{万} — three strokes: a long bar, a stroke that runs right and turns down with a hook, and a long stroke falling to the left. The next lesson puts it in your hand.

\begin{cousinweb}
\textbf{\zh{千万}} (\emph{qiānwàn}) is built from \zh{千} 'a thousand' and \zh{万} 'ten thousand', both inherited from Middle Chinese, \emph{tshen} and \emph{mjonH} (Guangyun, in Baxter's transcription), which glosses the first 'ten hundreds'; for the second it names a dance and an insect, and the number is not given there. A thousand times, ten thousand times over: the urging sense is usually explained that way.
\end{cousinweb}

\begin{grammarlens}[title={What you've built}]
One more word for saying what would happen.
\end{grammarlens}

\subsection*{Your turn}
\begin{itemize}
\raggedright
  \item \emph{Say it:} \emph{qiānwàn}
  \item \emph{Say it:} \emph{qiānwàn}, once more
  \item \emph{Say it:} \emph{nǐ yí ge rén zài wàibian qiānwàn yào bǎozhòng} — Out there on your own, be sure to look after yourself.
  \item \emph{Recall:} say the Mandarin for to change one's tune, then the Mandarin for to exhort (someone), then say \emph{qiānwàn} again
\end{itemize}

\begin{tcolorbox}[breakable,colback=teal!4,colframe=teal!35!black,title={Before you move on}]
What does \zh{千万} mean? (\textquotedblleft{}Be sure to (urging); whatever you do\textquotedblright{}.) Say it once more.
\end{tcolorbox}
\section[\texorpdfstring{\zh{万} (wàn)}{万 (wàn)}]{\zh{万} — a bar, a turning hook, a long fall}

\label{lesson:ZH-W393-wan-myriad}



\begin{quote}
The word first, then the character inside it.

\begin{itemize}
\raggedright
  \item \emph{Recall:} say \emph{qiānwàn} — \textbf{R1}, one lesson back, before it decays
\end{itemize}

\begin{quote}
\zh{千} \zh{为}
\end{quote}

\emph{Read it:} each one aloud, with nothing said first — \emph{qiān}, \emph{wéi}
\end{quote}

\subsection*{Script — a bar, a turning hook, a long fall}
\hlblockfigure{\detokenize{figures/ZH-W393-wan-myriad-filmstrip.pdf}}{How \zh{万} is written, stroke by stroke}

\begin{quote}
\zh{万}
\end{quote}

\textbf{Three strokes, top to bottom.}

\textbf{A long bar} across the top.

\textbf{A turning hook:} from under the middle of the bar, a stroke that runs right, turns down and hooks to the left, as in \zh{为} and \zh{月}.

\textbf{A long stroke falling to the left,} from the bar down to the lower left.

One, two, three.

\subsection*{Your turn}
\begin{itemize}
\raggedright
  \item \emph{Write it:} \zh{万}, all three, with the model in front of you
  \item \emph{Cover:} write it again from memory, and count the pieces, not the strokes
\end{itemize}

\begin{tcolorbox}[breakable,colback=teal!4,colframe=teal!35!black,title={Before you move on}]
How many strokes in \zh{万}? (\textbf{Three}.) Which stroke has the hook? (\textbf{The second, the turning hook}.) Which comes last? (\textbf{The long fall to the left}.)

Source: \href{https://raw.githubusercontent.com/chanind/hanzi-writer-data/68d10a4b21150cae5e1ebbd223eed289cf32d90c/data/\%E4\%B8\%87.json}{Hanzi Writer Data \zh{万}, pinned stroke paths and medians}.
\end{tcolorbox}
\section[\texorpdfstring{\zh{万能} (wànnéng)}{万能 (wànnéng)}]{\zh{万能} (wànnéng) — able to do everything; all-purpose}

\label{lesson:ZH-C393-wanneng}



\begin{quote}
Before the new one: say the Mandarin for to exhort (someone), then the Mandarin for be sure to (urging).
\end{quote}

\subsection*{You'll want to know: \zh{万能}}
\textbf{\zh{万能}} — \emph{wànnéng} — \textquotedblleft{}able to do everything; all-purpose\textquotedblright{}.

\begin{quote}
\emph{méiyǒu wànnéng de bànfǎ} — There's no way that works for everything.
\end{quote}

\begin{cousinweb}
\textbf{\zh{万能}} (\emph{wànnéng}) is built from \zh{万} 'ten thousand' and \zh{能} 'able', both inherited from Middle Chinese, \emph{mjonH} and \emph{nong} (Guangyun, in Baxter's transcription), which glosses the second 'skilled': able in ten thousand things.
\end{cousinweb}

\begin{grammarlens}[title={What you've built}]
One more word for saying what would happen.
\end{grammarlens}

\subsection*{Your turn}
\begin{itemize}
\raggedright
  \item \emph{Say it:} \emph{wànnéng}
  \item \emph{Say it:} \emph{wànnéng}, once more
  \item \emph{Say it:} \emph{méiyǒu wànnéng de bànfǎ} — There's no way that works for everything.
  \item \emph{Recall:} say the Mandarin for to exhort (someone), then the Mandarin for be sure to (urging), then say \emph{wànnéng} again
\end{itemize}

\begin{tcolorbox}[breakable,colback=teal!4,colframe=teal!35!black,title={Before you move on}]
What does \zh{万能} mean? (\textquotedblleft{}Able to do everything; all-purpose\textquotedblright{}.) Say it once more.
\end{tcolorbox}
\section[suíshí]{suíshí — at any time; whenever needed}

\label{lesson:ZH-C393-suishi}



\begin{quote}
Before the new one: say the Mandarin for be sure to (urging), then the Mandarin for all-purpose.
\end{quote}

\subsection*{You'll want to know: suíshí}
\emph{suíshí} — \textquotedblleft{}at any time; whenever needed\textquotedblright{}.

\begin{quote}
\emph{yàoshi yǒushì nǐ suíshí lái zhǎo wǒ} — If anything comes up, come and find me any time.
\end{quote}

Its characters come later; for now it is a word for the ear.

\begin{cousinweb}
\emph{suíshí} is built from \emph{suí} 'to follow' and \zh{时} 'the time', both inherited from Middle Chinese, \emph{zjwe} and \emph{dzyi} (Guangyun, in Baxter's transcription), which glosses the first 'to go along with': whatever the time.
\end{cousinweb}

\begin{grammarlens}[title={What you've built}]
One more word for saying what would happen.
\end{grammarlens}

\subsection*{Your turn}
\begin{itemize}
\raggedright
  \item \emph{Say it:} \emph{suíshí}
  \item \emph{Say it:} \emph{suíshí}, once more
  \item \emph{Say it:} \emph{yàoshi yǒushì nǐ suíshí lái zhǎo wǒ} — If anything comes up, come and find me any time.
  \item \emph{Recall:} say the Mandarin for be sure to (urging), then the Mandarin for all-purpose, then say \emph{suíshí} again
\end{itemize}

\begin{tcolorbox}[breakable,colback=teal!4,colframe=teal!35!black,title={Before you move on}]
What does suíshí mean? (\textquotedblleft{}At any time; whenever needed\textquotedblright{}.) Say it once more.
\end{tcolorbox}
\section[shùnlì]{shùnlì — smooth; without a hitch}

\label{lesson:ZH-C393-shunli}



\begin{quote}
Before the new one: say the Mandarin for all-purpose, then the Mandarin for at any time at all.
\end{quote}

\subsection*{You'll want to know: shùnlì}
\emph{shùnlì} — \textquotedblleft{}smooth; without a hitch\textquotedblright{}.

\begin{quote}
\emph{rúguǒ shùnlì dehuà wǒmen wǔ diǎn jiù dào} — If all goes well, we'll be there by five.
\end{quote}

Its characters come later; for now it is a word for the ear.

\begin{cousinweb}
\emph{shùnlì} is built from \emph{shùn} 'to go along with' and \emph{lì} 'favourable', both inherited from Middle Chinese, \emph{zywinH} and \emph{lijH} (Guangyun, in Baxter's transcription), which glosses the first 'to follow' and the second 'auspicious'.
\end{cousinweb}

\begin{grammarlens}[title={What you've built}]
One more word for saying what would happen.
\end{grammarlens}

\subsection*{Your turn}
\begin{itemize}
\raggedright
  \item \emph{Say it:} \emph{shùnlì}
  \item \emph{Say it:} \emph{shùnlì}, once more
  \item \emph{Say it:} \emph{rúguǒ shùnlì dehuà wǒmen wǔ diǎn jiù dào} — If all goes well, we'll be there by five.
  \item \emph{Recall:} say the Mandarin for all-purpose, then the Mandarin for at any time at all, then say \emph{shùnlì} again
\end{itemize}

\begin{tcolorbox}[breakable,colback=teal!4,colframe=teal!35!black,title={Before you move on}]
What does shùnlì mean? (\textquotedblleft{}Smooth; without a hitch\textquotedblright{}.) Say it once more.
\end{tcolorbox}
